\begin{frame}[fragile]{Comments (Hoe Maak Je Comments?)}
    In \LaTeX{} kun je -- net als in andere programmeertalen -- comments maken binnen je code. 
    Je geeft aan wat een comment moet zijn door een \%-teken aan het begin van de regel te zetten. 
    Je kunt in Overleaf ook \;\mintinline[bgcolor={gray!60!white}]{tex}{Ctrl + ?} \; typen.
    Dan komt er een \% aan het begin van de regel. \pause \LaTeX{} negeert \emph{alles} na het \%-teken. 
    Voorbeeld:
    \begin{minted}{tex}
        \documentclass[a4paper]{article}
        \begin{document}
        Hier is tekst die wel weergegeven wordt. 
        % Deze tekst wordt niet weergegeven. 
        \end{document}
    \end{minted}
    \pause 
    Dit geeft dan:
    \begin{tcolorbox}[boxrule=0.5pt, colback=white]
        Hier is tekst die wel weergegeven wordt. 
        % Deze tekst wordt niet weergegeven. 
    \end{tcolorbox}
\end{frame}
